\documentclass{article}
\usepackage{amsmath,amsthm,amssymb}
\usepackage{geometry}
\usepackage{mathrsfs}
\usepackage{enumitem}
\usepackage{graphicx}
\usepackage{hyperref}

\geometry{margin=1in}
\title{Rigorous Analysis: Bounding Box Algorithm Worst-Case Construction}
\author{}
\date{}

\newtheorem{definition}{Definition}
\newtheorem{lemma}{Lemma}
\newtheorem{theorem}{Theorem}
\newtheorem*{corollary}{Corollary}

\begin{document}

\maketitle

\section{Recursive Worst-Case Instance}

\begin{definition}[Recursive Point Set]
Let $n = 2^k$ for some integer $k \geq 1$. We define a point set $P_n \subset \mathbb{R}^2$ recursively as follows:
\begin{itemize}
    \item \textbf{Base case:} $P_2 = \{(0,0), (1,0)\}$
    \item \textbf{Recursive step:} For $n = 2^j$, define:
    \[
        P_n = T_L \cup T_R
    \]
    where:
    \begin{align*}
        T_L &= P_{n/2} \\
        T_R &= \{(x + D_n, y) : (x,y) \in P_{n/2} \} \\
        D_n &= 4n
    \end{align*}
\end{itemize}
\end{definition}

\paragraph{Geometric Properties}
\begin{itemize}
    \item All points lie along the line $y = 0$
    \item Total number of points: $|P_n| = n$
    \item Each recursive cluster $P_{n/2}$ lies in a contiguous horizontal segment of unit height and width $\approx n/2$
    \item At level $j$, the width of the entire point set is:
    \[
        \text{width}(P_n) = D_n + \text{diameter}(P_{n/2}) = \Theta(n)
    \]
\end{itemize}

\section{Behavior of the Bounding Box Algorithm}

We analyze the standard bounding box divide-and-conquer heuristic:
\begin{enumerate}
    \item Compute the bounding box of the point set
    \item If width $>$ height, split with a vertical line; otherwise, split horizontally
    \item Recursively solve left/right (or top/bottom) subproblems
    \item Connect the two paths by choosing the shortest among the four possible endpoint-to-endpoint connections
\end{enumerate}

\begin{lemma}[Splitting Direction]
At every level, the algorithm splits the instance $P_n$ with a vertical line, separating $T_L$ and $T_R$.
\end{lemma}

\begin{proof}
Since all points lie on $y = 0$, the bounding box has height $0$ and width $D_n + O(n) = \Theta(n)$. Thus, the algorithm always chooses a vertical split.
\end{proof}

\begin{lemma}[Length of Connecting Edge]
The minimum distance between any point in $T_L$ and any point in $T_R$ is exactly $D_n = 4n$.
\end{lemma}

\begin{proof}
All $x$-coordinates in $T_L$ lie in $[0, O(n)]$, while all $x$-coordinates in $T_R$ lie in $[4n, 4n + O(n)]$. Therefore, the shortest possible connection spans $D_n = 4n$ units.
\end{proof}

\section{Cost of the Algorithm's Output}

Let $W_A(n)$ denote the total Wiener index of the path produced by the bounding box algorithm on $P_n$. We define the recurrence:
\[
    W_A(n) = 2W_A(n/2) + C(n), \quad W_A(2) = 1
\]
where $C(n)$ is the contribution to the Wiener index from the connection between $T_L$ and $T_R$.

\begin{lemma}[Connection Cost]
\[
    C(n) = 4n \cdot \left(\frac{n}{2}\right)\left(\frac{n}{2}\right) = n^3
\]
\end{lemma}

\begin{proof}
\begin{itemize}
    \item The connecting edge has length $4n$
    \item In the final Hamiltonian path, this edge joins two subpaths of size $n/2$
    \item Its contribution to the Wiener index is:
    \[
        \text{length} \times \text{position weight} = 4n \cdot \frac{n}{2} \cdot \frac{n}{2} = n^3
    \]
\end{itemize}
\end{proof}

\begin{theorem}[Bounding Box Algorithm Cost]
\[
    W_A(n) = \Theta(n^3 \log n)
\]
\end{theorem}

\begin{proof}
We have the recurrence:
\[
    W_A(n) = 2W_A(n/2) + n^3
\]
Apply the Master Theorem:
\begin{itemize}
    \item $a = 2$, $b = 2$, $f(n) = n^3$
    \item $n^{\log_b a} = n$
    \item Since $f(n) = \omega(n^{\log_b a})$ and the regularity condition holds:
    \[
        af(n/b) = 2 \cdot (n/2)^3 = \frac{n^3}{4} < cf(n) \text{ for } c = \frac{1}{2} < 1
    \]
\end{itemize}
The recurrence falls under Case 3 of the Master Theorem. Therefore:
\[
    W_A(n) = \Theta(n^3 \log n)
\]
\end{proof}

\section{Optimal Solution Analysis}

Let $W^*(n)$ denote the Wiener index of the optimal Hamiltonian path on $P_n$.

\begin{theorem}[Optimal Cost]
\[
    W^*(n) = \Theta(n^3)
\]
\end{theorem}

\begin{proof}
Construct the path $\pi^*$ that visits:
\begin{itemize}
    \item All points in $T_L$ in increasing $x$-order
    \item A single jump from the rightmost point in $T_L$ to the leftmost point in $T_R$
    \item All points in $T_R$ in increasing $x$-order
\end{itemize}
\begin{itemize}
    \item All intra-cluster edges are of length $O(1)$ with total weight $O(n^2)$
    \item The single inter-cluster edge has length $4n$ and weight $\left(\frac{n}{2}\right)^2$
\end{itemize}
Therefore:
\[
    W^*(n) = 2 \cdot O(n) \cdot O(n^2) + 4n \cdot \left(\frac{n}{2}\right)^2 = O(n^3) + n^3 = \Theta(n^3)
\]
\end{proof}

\section{Approximation Ratio}

\begin{theorem}[Bounding Box Approximation Factor]
\[
    \frac{W_A(n)}{W^*(n)} = \Theta(\log n)
\]
\end{theorem}

\begin{proof}
From Theorems 1 and 2:
\[
    \frac{W_A(n)}{W^*(n)} = \frac{\Theta(n^3 \log n)}{\Theta(n^3)} = \Theta(\log n)
\]
\end{proof}

\begin{corollary}[Unbounded Approximation Ratio]
The bounding box algorithm does not admit a constant-factor approximation for the Hamiltonian path Wiener index problem in the Euclidean plane.
\end{corollary}

\end{document}

